%======================================================================
% a) Solucion por MEF
%======================================================================
\newpage
\section*{a) Soluci\'on por el m\'etodo de elementos finitos}

El problema se resolvi\'o con un c\'odigo propio de elasticidad plana en Python,
ordenado por etapas igual que en las tareas anteriores (malla $\to$
interpolaci\'on $\to$ matrices $\to$ resolver $\to$ graficar). El c\'odigo est\'a
en el anexo.

%----------------------------------------------------------------------
\begin{desarrollo}{1. Par\'ametros del problema}

El enunciado deja todos los par\'ametros libres. Se eligi\'o una placa de acero
grande respecto al agujero ($L=20a$), para que la soluci\'on de placa infinita
sea una buena referencia.

\begin{center}
\renewcommand{\arraystretch}{1.2}
\begin{tabular}{@{}lllc@{}}
\toprule
& Par\'ametro & S\'imbolo & Valor \\
\midrule
\multirow{5}{*}{Geometr\'ia}
& Radio del agujero & $r=a$ & $10$ mm \\
& Di\'ametro del agujero & $2r$ & $20$ mm \\
& Semiancho = semilargo de la placa & $L$ & $200$ mm \\
& Ancho total & $d=2L$ & $400$ mm \\
& Espesor & $t$ & $10$ mm \\
\midrule
\multirow{3}{*}{Material (acero)}
& M\'odulo de Young & $E$ & $200$ GPa \\
& Coeficiente de Poisson & $\nu$ & $0.3$ \\
& Estado & -- & tensi\'on plana \\
\midrule
\multirow{2}{*}{Carga}
& Tracci\'on en $x=\pm L$ & $S$ & $100$ MPa \\
& Resultante en el 1/4 de placa & $S\,t\,L$ & $200$ kN \\
\midrule
\multirow{2}{*}{Relaciones}
& Ancho / di\'ametro & $d/2r$ & $20$ \\
& Largo cargado / radio & $L/a$ & $20$ \\
\bottomrule
\end{tabular}
\end{center}

\begin{center}
\renewcommand{\arraystretch}{1.2}
\begin{tabular}{@{}ll@{}}
\toprule
Par\'ametro de la discretizaci\'on & Valor \\
\midrule
Dominio modelado & 1/4 de placa ($x\ge0$, $y\ge0$), por simetr\'ia \\
Tipo de elemento & Q9 isoparam\'etrico (9 nodos, 18 GDL) \\
Mallas & malla 1: 80 Q9 (714 GDL); malla 2: 320 Q9 (2706 GDL) \\
Graduaci\'on radial & raz\'on 15 entre el \'ultimo y el primer elemento \\
Integraci\'on en el interior ($\mathbf{K}^{(e)}$) & Gauss--Legendre $3\times3$ (9 puntos) \\
Integraci\'on en los bordes ($\mathbf{f}^{(e)}$) & Gauss--Legendre 1D, 4 puntos por lado \\
Tensiones & en los nodos, promediadas entre elementos vecinos \\
\bottomrule
\end{tabular}
\end{center}

\textbf{Justificaci\'on breve.}
\begin{itemize}[itemsep=2pt]
  \item $d/2r=20\ge4$: seg\'un Timoshenko, con $d/2r\ge4$ el error de usar la
  soluci\'on de placa infinita para $(\sigma_\theta)_{\max}$ ya es menor al 6\,\%.
  \item $L/a=20$: el borde cargado queda lejos del agujero. Con una placa m\'as
  chica ($L=5a$) se obtuvo $K_t\approx3.36$, un 12\,\% sobre Kirsch, por el
  tama\~no finito de la placa; con $L=20a$ esa diferencia baja a menos del 1\,\%
  (ver b).
  \item Tensi\'on plana: la placa es delgada y sin cargas fuera del plano.
  \item 1/4 de placa: la geometr\'ia y la carga son sim\'etricas respecto a $x$ e $y$.
\end{itemize}
\end{desarrollo}

%----------------------------------------------------------------------
\begin{desarrollo}{2. Modelo y condiciones de borde}

\begin{center}
\begin{tikzpicture}[>=latex, scale=0.9]
    \def\Lq{5}   % L
    \def\aq{1}   % a
    % cuarto de placa
    \draw[thick, fill=black!5]
        (\aq,0) -- (\Lq,0) -- (\Lq,\Lq) -- (0,\Lq) -- (0,\aq)
        arc[start angle=90, end angle=0, radius=\aq] -- cycle;
    % apoyos deslizantes en y = 0 (u_y = 0)
    \foreach \x in {1.4,2.2,...,4.8}{
        \draw (\x,-0.13) circle (0.12);
    }
    \draw[thick] (\aq-0.1,-0.27) -- (\Lq+0.1,-0.27);
    \foreach \x in {1.0,1.3,...,5.0}{\draw (\x,-0.27) -- (\x-0.15,-0.45);}
    % apoyos deslizantes en x = 0 (u_x = 0)
    \foreach \y in {1.4,2.2,...,4.8}{
        \draw (-0.13,\y) circle (0.12);
    }
    \draw[thick] (-0.27,\aq-0.1) -- (-0.27,\Lq+0.1);
    \foreach \y in {1.0,1.3,...,5.0}{\draw (-0.27,\y) -- (-0.45,\y-0.15);}
    % carga S en x = L
    \foreach \y in {0.25,0.75,...,4.75}{
        \draw[-{Latex[length=2mm]}, thick, red!70!black] (\Lq,\y) -- (\Lq+0.8,\y);
    }
    \node[red!70!black] at (\Lq+1.25,\Lq/2) {\large $S$};
    % ejes
    \draw[-{Latex[length=2mm]}, gray] (0,0) -- (1.9,0) node[above right=-1pt] {$x$};
    \draw[-{Latex[length=2mm]}, gray] (0,0) -- (0,1.9) node[above right=-1pt] {$y$};
    % puntos de control
    \foreach \p/\pos/\lab in {{(\aq,0)}/above right/A, {(0,\aq)}/above right/B,
                              {(\Lq,0)}/above left/C, {(0,\Lq)}/below right/D}{
        \fill[ucred] \p circle (2pt) node[\pos=1pt, text=ucred] {\small\lab};
    }
    % cotas
    \draw[<->, thin] (0,0) -- ({\aq*cos(45)},{\aq*sin(45)});
    \node at (0.25,0.55) {\small $a$};
    \draw[<->, thin] (0,-0.8) -- (\Lq,-0.8) node[midway, fill=white, inner sep=1.5pt] {\small $L$};
    \draw[<->, thin] (-0.8,0) -- (-0.8,\Lq) node[midway, fill=white, inner sep=1.5pt] {\small $L$};
    % etiquetas de borde
    \node[font=\scriptsize, align=center] at (2.6,\Lq+0.3) {borde libre};
    \node[font=\scriptsize, align=center, rotate=-45] at (1.05,1.05) {libre};
\end{tikzpicture}

{\small Esquema, no a escala: en el modelo $L=20a$.}

\end{center}


\begin{center}
\renewcommand{\arraystretch}{1.2}
\begin{tabular}{@{}lll@{}}
\toprule
Borde & Condici\'on & C\'omo se impuso \\
\midrule
$y=0$ (simetr\'ia) & $u_y=0$, $u_x$ libre & se eliminan los GDL $u_y$ de sus nodos \\
$x=0$ (simetr\'ia) & $u_x=0$, $u_y$ libre & se eliminan los GDL $u_x$ de sus nodos \\
$x=L$ & $\bar{\mathbf{t}}=(S,0)$ & fuerzas nodales consistentes (ver punto 3) \\
$y=L$ y $r=a$ & $\bar{\mathbf{t}}=\mathbf{0}$ (libre) & nada (no aportan a $\mathbf{F}$) \\
\bottomrule
\end{tabular}
\end{center}

Las dos condiciones de simetr\'ia eliminan tambi\'en los 3 movimientos de cuerpo
r\'igido, as\'i que no hace falta otro apoyo. Los puntos de control para comparar
con Kirsch son \textbf{A} $=(a,0)$ y \textbf{B} $=(0,a)$ en el borde del agujero,
y \textbf{C} $=(L,0)$ y \textbf{D} $=(0,L)$ en el contorno exterior.
\end{desarrollo}

%----------------------------------------------------------------------
\begin{desarrollo}{3. Integraci\'on num\'erica}

Con la formulaci\'on de desplazamientos se llega a $\mathbf{K}\,\mathbf{U}=\mathbf{F}$,
con dos tipos de integrales por elemento: una en el \textbf{interior} (rigidez) y
otra en los \textbf{lados cargados} (fuerzas). Ambas se calculan con Gauss--Legendre
en el elemento de referencia.

\textbf{Interior: matriz de rigidez.}
\[
\mathbf{K}^{(e)}=\int_{A^{(e)}}\mathbf{B}^T\mathbf{D}\,\mathbf{B}\,t\,dA
\;\approx\;\sum_{i=1}^{9} w_i\,\mathbf{B}_i^T\mathbf{D}\,\mathbf{B}_i\,t\,\det\mathbf{J}_i
\]
\begin{itemize}[itemsep=2pt]
  \item Regla $3\times3$: $\xi,\eta\in\{-\sqrt{3/5},0,\sqrt{3/5}\}$, pesos
  $\{\tfrac{25}{81},\tfrac{40}{81},\tfrac{64}{81}\}$.
  \item Es integraci\'on completa para el Q9. 
  \item En cada punto se eval\'ua $\mathbf{J}$, y si $\det\mathbf{J}\le0$ el
  c\'odigo se detiene (elemento invertido o muy distorsionado).
  \item Los elementos junto al agujero son curvos, as\'i que la regla no es
  exacta ah\'i, pero su error es despreciable frente al de la malla.
\end{itemize}

\textbf{Bordes: fuerzas nodales equivalentes.}
\[
\mathbf{f}_i=\int_{\Gamma^{(e)}}N_i\,\mathbf{q}\,ds
\;\approx\;\sum_{p=1}^{4} w_p\,N_i(s_p)\,\mathbf{q}(s_p)\left\lVert\tfrac{d\mathbf{x}}{ds}\right\rVert_p,
\qquad \mathbf{q}=\bar{\mathbf{t}}\,t
\]
\begin{itemize}[itemsep=2pt]
  \item Cada lado del Q9 tiene 3 nodos y se parametriza con $s\in[0,1]$.
  \item Se usan 4 puntos de Gauss, exacto en lados rectos (integrando de grado 3).
  \item Con $q$ constante el reparto es $q\ell\,(\tfrac16,\tfrac46,\tfrac16)$:
  el nodo del medio se lleva 4 veces lo de cada esquina.
  \item Control: $\sum F_x = S\,t\,L = 200$ kN, que es lo que entrega el c\'odigo
  en las dos mallas.
\end{itemize}

\begin{center}
\figuraT{T3/figuras/fuerzas_nodales_problema_2.png}{0.7}{Fuerzas nodales equivalentes de $S$ en $x=L$.}
\end{center}
\end{desarrollo}

%----------------------------------------------------------------------
\begin{desarrollo}{4. Mallado}

\begin{itemize}[itemsep=2pt]
  \item El cuarto de placa se divide en dos bloques separados por la diagonal
  $y=x$: arco $0^\circ$--$45^\circ$ hasta $x=L$, y arco $45^\circ$--$90^\circ$
  hasta $y=L$.
  \item Los elementos crecen con una progresion generando de esta manera que los elementos mas pequeños esten cerca de la zona de concentracion de esfuerzos.
  Como la placa es grande ($L=20a$), se usa una raz\'on 15 entre el \'ultimo y
  el primer elemento radial. Sobre el eje $x$, en la malla 2 el primer elemento
  mide 1.8 mm y el \'ultimo 27 mm (en la malla 1, 3.5 y 52 mm): chicos donde
  est\'a el gradiente y grandes donde el campo ya es casi uniforme.
  \item Los elementos del agujero tienen el lado curvo (par\'abola por 3 puntos
  del arco), as\'i que el contorno queda bien representado.
\end{itemize}

\begin{center}
\renewcommand{\arraystretch}{1.2}
\begin{tabular}{@{}lcccc@{}}
\toprule
Malla & \shortstack{divisiones por bloque\\(radial $\times$ angular)} & elementos Q9 & nodos & GDL\\
\midrule
Malla 1 (\texttt{problema\_2}) & $10\times4$ & 80 & 357 & 714\\
Malla 2 (\texttt{problema\_3}) & $20\times8$ & 320 & 1353 & 2706\\
\bottomrule
\end{tabular}
\end{center}

\begin{center}
\figuraT{T3/figuras/malla_problema_2.png}{0.48}{Malla 1: 80 Q9, con apoyos y carga.}\hfill
\figuraT{T3/figuras/malla_problema_3.png}{0.48}{Malla 2: 320 Q9, con apoyos y carga.}
\end{center}


Para este caso se comparan ambos casos para asegurar que no sea dependiente de la malla y que sea el resultado sea representativo de la carga

\end{desarrollo}

%----------------------------------------------------------------------
\begin{desarrollo}{5. Resoluci\'on y c\'alculo de tensiones}

\begin{itemize}[itemsep=2pt]
  \item \textbf{Ensamble}: $\mathbf{K}^{(e)}$ ($18\times18$) se suma en los GDL
  $[2n_k,\,2n_k+1]$ de sus 9 nodos.
  \item \textbf{Resoluci\'on}: se eliminan los GDL restringidos, se resuelve
  $\mathbf{K}_{red}\mathbf{U}_{red}=\mathbf{F}_{red}$ y se reconstruye $\mathbf{U}$
  (igual que en T1 y T2).
  \item \textbf{Tensiones}: $\bm\sigma=\mathbf{D}\,\mathbf{B}\,\mathbf{u}^{(e)}$
  evaluada en los 9 nodos de cada elemento. Como $\bm\sigma$ salta entre
  elementos, en cada nodo se promedian los valores de los elementos vecinos.
  \item \textbf{Von Mises}:
  $\sigma_{VM}=\sqrt{\sigma_{xx}^2-\sigma_{xx}\sigma_{yy}+\sigma_{yy}^2+3\tau_{xy}^2}$.
\end{itemize}
\end{desarrollo}

%----------------------------------------------------------------------
\begin{desarrollo}{6. Resultados del MEF}

\begin{center}
\figuraT{T3/figuras/tensiones_problema_2.png}{0.8}{Campos $\sigma_{xx}$, $\sigma_{yy}$, $\tau_{xy}$ y von Mises con la malla 1, en la zona del agujero ($0\le x,y\le5a$; m\'as lejos el campo es casi uniforme).}

\vspace{6pt}
\figuraT{T3/figuras/tensiones_problema_3.png}{0.8}{Campos $\sigma_{xx}$, $\sigma_{yy}$, $\tau_{xy}$ y von Mises con la malla 2, en la zona del agujero ($0\le x,y\le5a$).}

\vspace{6pt}
\figuraT{T3/figuras/deformada_mef.png}{0.6}{Deformada amplificada y desplazamientos $u_x$, $u_y$.}
\end{center}

\begin{center}
\renewcommand{\arraystretch}{1.25}
\begin{tabular}{@{}lcccccc@{}}
\toprule
& \multicolumn{3}{c}{Malla 1 (80 Q9)} & \multicolumn{3}{c}{Malla 2 (320 Q9)}\\
\cmidrule(lr){2-4}\cmidrule(lr){5-7}
Punto & $u$ [mm] & $\sigma_{xx}$ [MPa] & $\sigma_{yy}$ [MPa] & $u$ [mm] & $\sigma_{xx}$ [MPa] & $\sigma_{yy}$ [MPa]\\
\midrule
A $(a,0)$ & $u_x=0.01510$  & $-10.5$ & $-104.9$ & $u_x=0.01511$  & $-4.2$  & $-103.0$\\
B $(0,a)$ & $u_y=-0.00508$ & $307.2$ & $16.4$   & $u_y=-0.00508$ & $304.0$ & $6.1$\\
C $(L,0)$ & $u_x=0.10126$  & $100.1$ & $1.7$    & $u_x=0.10127$  & $100.0$ & $1.7$\\
D $(0,L)$ & $u_y=-0.03075$ & $99.1$  & $-0.1$   & $u_y=-0.03076$ & $99.1$  & $0.0$\\
\midrule
\multicolumn{7}{@{}l}{$K_t=\sigma_{xx}(B)/S$:\quad malla 1 = $3.07$ \qquad malla 2 = $3.04$}\\
\bottomrule
\end{tabular}
\end{center}-

\end{desarrollo}
